% Figure: recombinant insulin manufacture process flow:
% E. coli fermentation -> cell lysis -> IEX -> refolding -> polishing -> formulation.
\begin{figure}[htbp]
\centering
\definecolor{engorange}{RGB}{217, 119, 6}
\begin{tikzpicture}[
  block/.style={draw=engblack, thick, rounded corners,
                fill=engblue!10, minimum width=2.6cm, minimum height=1.0cm,
                font=\small, align=center},
  fermblock/.style={block, fill=englime!20},
  recblock/.style={block, fill=engorange!15},
  refoldblock/.style={block, fill=engred!12},
  polishblock/.style={block, fill=engblue!15},
  formblock/.style={block, fill=engblue!8},
  arrow/.style={-{Stealth}, thick},
  lab/.style={font=\scriptsize, anchor=west, engblack!70, align=left},
]

\node[fermblock] (strain) at (0, 0) {Production\\strain (\emph{E.\ coli})};
\node[fermblock] (ferm) at (3.5, 0) {Fed-batch\\fermentation};
\node[fermblock] (harv) at (7.0, 0) {Cell harvest\\(centrifuge)};
\node[recblock] (lysis) at (10.5, 0) {Cell lysis\\(homogeniser)};

\draw[arrow] (strain) -- (ferm);
\draw[arrow] (ferm) -- (harv);
\draw[arrow] (harv) -- (lysis);

\draw[arrow] (lysis.south) -- ++(0, -0.5) -| (0, -1.8);

\node[recblock] (ibwash) at (0, -2.4) {Inclusion-body\\wash};
\node[recblock] (solubilise) at (3.5, -2.4) {Solubilisation\\(denaturant + reducing)};
\node[refoldblock] (refold) at (7.0, -2.4) {Oxidative\\refolding};
\node[refoldblock] (cleave) at (10.5, -2.4) {Enzymatic\\cleavage};

\draw[arrow] (ibwash) -- (solubilise);
\draw[arrow] (solubilise) -- (refold);
\draw[arrow] (refold) -- (cleave);

\draw[arrow] (cleave.south) -- ++(0, -0.5) -| (0, -4.2);

\node[polishblock] (iex1) at (0, -4.8) {IEX\\chromatography};
\node[polishblock] (rpc) at (3.5, -4.8) {Reversed-phase\\chromatography};
\node[polishblock] (gel) at (7.0, -4.8) {Size-exclusion\\polishing};
\node[formblock] (cryst) at (10.5, -4.8) {Crystallisation\\(Zn-insulin)};

\draw[arrow] (iex1) -- (rpc);
\draw[arrow] (rpc) -- (gel);
\draw[arrow] (gel) -- (cryst);

\draw[arrow] (cryst.south) -- ++(0, -0.5) -| (5.0, -6.0);

\node[formblock] (formulate) at (5.0, -6.5) {Formulation\\(buffer, Zn, preservative)};
\node[formblock] (fill) at (9.0, -6.5) {Sterile fill\\(vials, cartridges)};

\draw[arrow] (formulate) -- (fill);

% Row labels
\node[lab] at (-1.7, 0) {\textbf{Upstream}};
\node[lab] at (-1.7, -2.4) {\textbf{Recovery}};
\node[lab] at (-1.7, -4.8) {\textbf{Polishing}};
\node[lab] at (-1.7, -6.5) {\textbf{Finish}};

\end{tikzpicture}
\caption{Recombinant human insulin manufacture (toy schematic). A
production \emph{E.\ coli} strain carrying a chemically synthesised
human-insulin sequence is grown in a fed-batch fermenter (typical
working volumes 10000 to 50000 L), and the product expresses as
inclusion bodies. The cells are harvested, the inclusion bodies
solubilised in denaturant plus reducing agent, and the linear chain is
oxidatively refolded with three native disulfides at low protein
concentration. Enzymatic cleavage (for a proinsulin or fusion-protein
route) liberates the mature two-chain hormone. Ion-exchange,
reversed-phase, and size-exclusion chromatography polish the
preparation. Zinc-insulin crystallisation gives a stable bulk
intermediate; formulation, sterile fill, and finish complete the
process. The Eli Lilly Humulin process (approved 1982 in the United
States) was the first commercial recombinant biopharmaceutical and
followed approximately this engineering structure with insulin
A-chain and B-chain produced separately, then joined chemically,
under license from Genentech (\cite{paper:v6c10-goeddel-insulin-1979},
\cite{paper:v6c10-humulin-eli-lilly-1982}).}
\label{fig:vol06ch10:insulin-process}
\end{figure}
